% Generated draft: requires mathematical and source review.
\documentclass[11pt]{article}
\usepackage[margin=25mm]{geometry}
\usepackage{fontspec}
\setmainfont{Latin Modern Roman}
\usepackage{amsmath,amssymb,mathtools,unicode-math}
\setmathfont{Latin Modern Math}
\usepackage{enumitem,needspace,xurl,hyperref}
\setlength{\parindent}{0pt}
\setlength{\parskip}{5pt}
\setlength{\emergencystretch}{3em}
\newcommand{\N}{\mathbb{N}}
\newcommand{\Z}{\mathbb{Z}}
\newcommand{\Q}{\mathbb{Q}}
\newcommand{\R}{\mathbb{R}}
\newcommand{\C}{\mathbb{C}}
\newcommand{\F}{\mathbb{F}}
\newcommand{\A}{\mathbb{A}}
\newcommand{\PP}{\mathbb{P}}
\newcommand{\E}{\mathbb{E}}
\newcommand{\eps}{\varepsilon}
\newcommand{\dd}{\,\mathrm{d}}
\begin{document}
\section*{UnsolvedMath ID: 11000210}
\textbf{Source problem code:} catalogue entry
\subsection{Definitions and mathematical statement}
For a planar polygon $P$, its billiard dynamics moves a unit-speed point along straight segments and reflects at an edge by equal incoming and outgoing angles. If $n$ is a unit normal at an edge, the velocity changes by
\[
 v\longmapsto v-2(v\cdot n)n.
\]
Trajectories hitting vertices are singular and require exclusion or a specified continuation convention. A billiard system includes its polygon and the corresponding flow or collision map; directional restrictions may also be part of the data.

\textbf{Problem 3 (renormalization of billiards in polygons).} Is there a natural dynamical system on a space of polygonal billiard systems that yields a useful renormalization procedure? The desired dynamics acts on the space of billiard systems, not merely by moving a particle in a single fixed polygon. It should relate billiard dynamics at different scales or through an appropriate induction/rescaling scheme in a way useful for studying the original dynamics.

The excerpt does not define the parameter space, topology, equivalence under similarities, polygon class, directional data, or the precise relation between a renormalized billiard and the initial system. It also does not specify a mathematical criterion for ``natural'' or ``useful.'' In particular, it does not restrict to rational-angle polygons, convex polygons, or polygons unfolding to compact translation surfaces. No such restriction, or extension of a known rational-polygon framework to arbitrary polygons, is silently assumed. A precise existence theorem requires these choices and a stated dynamical conclusion to be supplied.
\subsection{Short English statement}
Construct a useful renormalization dynamics on the space of polygonal billiard systems. The source leaves the polygon class, parameter space, rescaling relation, and success criterion open.
\subsection{Sources}
\begin{itemize}
\item Local UnsolvedMath registry, permanent ID 11000210, source code catalogue entry; no proposer or proposal year supplied.
\item Farb, edited volume \emph{Problems on Mapping Class Groups and Related Topics}, Problem 3, PDF page 256; \url{https://www.math.uchicago.edu/~farb/papers/mcgbook.pdf}. The registry records PDF-font-structure extraction and access on 2026-07-29; its inventory distinguishes 329 labeled headings by source page.
\item No separate authoritative reference or candidate renormalization construction is supplied. The research summary says its literature status cannot be verified.
\end{itemize}
\end{document}
